\section{Scrapping (Vincent)}\label{sec:scrapping}

Scrapping decides how long old plants survive: it sets the spread of vintages behind
unequal margins, the speed at which a new technique spreads, and which plant sets the floor
of a price war. Three decisions are kept apart in the code: \emph{run or idle} (production,
in merit order), \emph{retire} (age or obsolescence, identical for both types) and
\emph{replace} (a second choice of technique, where the types differ through $h$).

\paragraph{The trap: replacement is a second choice of technique.} With the hurdle $h$,
replacing plant $v$ by technique $n$ pays when
\[
ucs_v > ucs_n + (\delta + h)\,\kappa_n .
\]
Shaikh firms ($h = 0$) replace earlier, which gives them younger capital and faster
diffusion; Okishio firms keep old plants longer. Retirement is therefore identical for both
types, and only replacement uses $h$, the same $h$ as investment.

\paragraph{Starting rule.} A plant retires at age $T = 1/\delta = 10$, or once its owner's
price has stayed below its $ucs$ for $N = 4$ periods (its quasi-rent is negative, Salter
1966). Replacement is on in version 1: the firm compares each working plant with the
technique its own hurdle picks, and replaces it if the inequality holds and the firm can pay
(cash, then debt up to the cap). Scrap value is zero; remaining book value is written off.
Depreciation is accounting only. When a firm's last plant retires, its lifetime profit rate
decides non-renewal (section~\ref{sec:entry}).

\paragraph{Tests} (\code{tests/test\_scrapping.py}): a plant reaching age $T$ leaves; a plant
whose $ucs$ stays above the price for $N$ periods leaves; an old plant with $ucs = 80$
facing a technique with $ucs = 60$, $\kappa = 150$: a Shaikh firm replaces it
($75 < 80$), an Okishio firm at $h = 16\,\%$ keeps it ($99 > 80$); on the test industry no
working plant is replaced; capacity after scrapping equals capacity before minus retired
plants.

\subsection*{Code}
\eqfile{retirement.py}
\eqfile{replacement.py}
\eqfile{lifetime_profit_rate.py}
\schedulerpart{scrapping}{\code{model/scheduler.py}: step 4}
